\documentclass[11pt]{article}
\usepackage[letterpaper, top=0.6in, bottom=0.6in, left=1in, right=1in, headheight=14pt]{geometry}
\usepackage[T1]{fontenc}
\usepackage{lmodern}
\usepackage{microtype}
\usepackage{enumitem}
\usepackage{titlesec}
\usepackage{tabularx}
\usepackage{array}
\usepackage{caption}
\usepackage{hhline}
\usepackage{fancyhdr}
\usepackage{lastpage}
\usepackage[most]{tcolorbox}
\usepackage[hidelinks, pdftitle={Open Orders Report Macro}, pdfauthor={Gust F.}]{hyperref}

\linespread{1.02}

\pagestyle{fancy}
\fancyhf{}
\fancyhead[L]{\small Strategy}
\fancyhead[R]{\small Rev.~12 \quad 2026-10-15}
\fancyfoot[L]{\small Gust F.}
\fancyfoot[C]{\small Page \thepage\ of \pageref{LastPage}}
\renewcommand{\headrulewidth}{0.4pt}

\titleformat{\section}{\normalsize\bfseries}{\thesection.}{0.6em}{}
\titlespacing*{\section}{0pt}{0.6em}{0.3em}
\setlist[enumerate]{leftmargin=2.2em, itemsep=0.1em, topsep=0.15em, label=(\alph*), ref=\alph*}
\setlist[itemize]{leftmargin=2.2em, itemsep=0.1em, topsep=0.15em, label=--}
\captionsetup[table]{font=small, labelfont=bf, labelsep=period, skip=0.4em, position=top, hypcap=false}
\setlength{\parindent}{0pt}
\setlength{\parskip}{0pt}

\newcommand{\menu}[1]{\textit{#1}}
\newcommand{\gt}{~>\ }
\renewcommand{\tabularxcolumn}[1]{m{#1}}
\newtcolorbox{notebox}{
  enhanced, breakable=false,
  colback=black!4, colframe=black, boxrule=0.8pt,
  left=5pt, right=5pt, top=2pt, bottom=2pt,
  before skip=3pt, after skip=3pt
}
\newcommand{\warnsign}{%
  \begin{tikzpicture}[baseline=-0.6ex]
    \draw[line width=0.8pt, fill=yellow!85, rounded corners=1pt]
      (0,0) -- (0.5,0) -- (0.25,0.43) -- cycle;
    \node[font=\bfseries\scriptsize] at (0.25,0.16) {!};
  \end{tikzpicture}%
}
\newcolumntype{L}{>{\raggedright\arraybackslash}X}

\begin{document}

\begin{center}
  {\Large\bfseries Open Orders Report Macro}\\[0.35em]
  {\normalsize User manual and road map}
\end{center}

\vspace{-0.4em}
\noindent\rule{\linewidth}{0.8pt}

\vspace{0.2em}
\textbf{Scope.} The Excel macro (v2.8) that builds \textit{001.\ Open Orders -- Order Intake -- Inventory Analysis By MSA}. It reads the Order Intake, Inventory by Month and Open Order exports for all MSAs from one folder and the size class workbook, and builds one print-ready workbook with banded tables: File Format, Order Intake (all MSAs), Inventory by Month (units by MSA, and the detail), Open Orders by MSA (open orders by MSA and type, charted) and Run Log.

\vspace{-0.2em}
\noindent\rule{\linewidth}{0.4pt}

\section{Key settings}\label{sec:settings}
\begin{center}
\small
\captionof{table}{Config settings to check before each run; leave all others at their defaults.}\label{tab:config}
\renewcommand{\arraystretch}{1.0}
\begin{tabularx}{\linewidth}{|>{\footnotesize\ttfamily\raggedright\arraybackslash}m{1.6in}|L|}
  \hline
  \multicolumn{1}{|l|}{\bfseries Setting} & \bfseries What it controls \\
  \hhline{|=|=|}
  InputFolder & Folder with the exports of all MSAs; blank opens a folder picker. \\
  \hline
  SizeClassFile & The size class workbook (models; classes on its Set-up tab). \\
  \hline
  MSAColumn & MSA column of the exports; blank finds it automatically. \\
  \hline
  InventoryMonth & Any day in the month to count; blank uses the latest month. \\
  \hline
  ShowCancelled & FALSE hides cancelled orders behind a filter. \\
  \hline
  ShowSSL & FALSE hides skid-steer loaders behind a filter and off the chart. \\
  \hline
\end{tabularx}
\end{center}

\section{One-time setup}\label{sec:setup}
\begin{notebox}
  \warnsign\hspace{0.5em}\textbf{Note:} You must save the macro in the blank worksheet itself, not in the personal worksheet.
\end{notebox}
\begin{enumerate}
  \item Save a blank Excel workbook as \texttt{Open Orders Report Builder.xlsm}.
  \item Alt+F11: import \texttt{build-open-orders-report.bas}, \menu{Debug\gt Compile VBAProject}, save.
  \item Run \texttt{ShowOpenOrdersReportConfig} (Alt+F8), the only macro listed, to add Config.
  \item Click \menu{Choose size class file...} on Config and pick the size class workbook.
  \item Click \menu{Enable Content} whenever Excel shows the macro security bar.
\end{enumerate}

\section{Running the report}\label{sec:run}
\begin{enumerate}
  \item Export the three reports for all MSAs into one folder.
  \item Check Table~\ref{tab:config}, then click \menu{Build report} on Config, the only way to build the report.
  \item Read the Run Log, Last run box and chart; a model the size class workbook does not list is UNMAPPED: add it there and rebuild.
  \item Filter Order Intake (e.g.\ MSA, then size class): its top row totals the visible rows.
  \item To see cancelled or skid-steer rows, clear the filter on that column.
  \item Print with Ctrl+P: each tab prints landscape, one page wide, with page numbers.
\end{enumerate}

\section{Road map}\label{sec:roadmap}
\begin{center}
\small
\captionof{table}{Road map from the current macro to automated PowerPoint charts.}\label{tab:roadmap}
\renewcommand{\arraystretch}{1.0}
\begin{tabularx}{\linewidth}{|>{\centering\arraybackslash}m{0.25in}|>{\raggedright\arraybackslash}m{0.95in}|L|>{\raggedright\arraybackslash}m{1.45in}|}
  \hline
  \bfseries No. & \bfseries Step & \bfseries Main actions & \bfseries Done when \\
  \hhline{|=|=|=|=|}
  1 & Confirm rules & Agreed 2026-10-12 and after; one item still open. & Last item answered. \\
  \hline
  2 & Build v2 & All MSAs, size list, hidden rows, style, chart, total. & Done 2026-10-15 (v2.8). \\
  \hline
  3 & Test in Excel & Run the test protocol against the hand-built tab. & The protocol passes. \\
  \hline
  4 & Adopt the list & The owner's size class list replaces the starter list. & The workbook is the only source of a class; nothing UNMAPPED. \\
  \hline
  5 & Feed the deck & An \menu{Update PowerPoint} button fills the deck charts. & Deck figures match. \\
  \hline
  6 & Hand over & Name owners; keep all files in one shared folder. & Someone else runs it. \\
  \hline
\end{tabularx}
\end{center}

\end{document}
